\documentclass[10pt,a4paper]{amsart}
\usepackage[margin=19mm]{geometry}
\usepackage{amsmath,amssymb,mathtools}
\usepackage[scr=rsfs,cal=euler]{mathalfa}
\usepackage{xfrac}
\usepackage[unicode,hidelinks,bookmarks=false]{hyperref}
\newcommand{\ie}{i.e.\ }
\newcommand{\cf}{cf.\ }
\newcommand{\resp}{respectively\ }
\newcommand{\Subsection}{\subsection}
\providecommand{\nobreakdash}{\nobreak\hbox{-}}
\setlength{\emergencystretch}{2em}
\multlinegap0pt
\usepackage[shortlabels]{enumitem}
\usepackage{amssymb}
\usepackage{mathtools}
\usepackage{float}
\usepackage{xcolor}
\newtheorem{as}{Assumption}
\newtheorem{lemma}{Lemma}
\newcommand{\Dfam}{\mathfrak{D}}     % scaling family
\newcommand{\Dm}{\mathcal{D}}        % scaling matrix
\newcommand{\R}{\mathbb{R}}
\newcommand{\Sol}{S}                 % solution set
\newcommand{\inner}[2]{\langle #1,\, #2\rangle}
\newcommand{\norm}[1]{\lVert #1 \rVert}
\title{\large A Variable-Metric Non-monotone Line Search Method for Mixed Variational Inequalities and Equilibrium Problems}
\author{Mohammed Alshahrani}
\date{Source: arXiv:2606.22869v1 (CC BY 4.0)}
\begin{document}
\maketitle

\noindent\emph{Source and licence.} \large A Variable-Metric Non-monotone Line Search Method for Mixed Variational Inequalities and Equilibrium Problems. Version arXiv:2606.22869v1 (CC BY 4.0), distributed by arXiv under the Creative Commons Attribution 4.0 International licence, http://creativecommons.org/licenses/by/4.0/, as recorded in the arXiv licence metadata for this exact version and shown on the abstract page https://arxiv.org/abs/2606.22869v1. Full licence notice: \texttt{licenses/arxiv-2606.22869-CC-BY-4.0.txt}.\par\smallskip

\noindent\textit{Editorial scope.} Counted unit: the article's lemma lem:descent (source0.tex lines 781--794). The article's complete original proof of it is delivered with it (source0.tex lines 795--837, one \texttt{proof} environment of 43 lines). No mathematics is written, repaired or re-derived: the statement and the proof are the article's own lines, in the article's own notation, and no retained line is substituted or re-flowed.  Retained uncounted as the source-local context the proof needs: lines 690--692; lines 694--703; lines 705--708; lines 719--722; lines 748--750.  Nothing was omitted from any retained span, so the package carries no omission record.  No retained line contains a citation, so the wrapper carries no bibliography and no bibliography entry.  No retained line contains a figure environment, an image-inclusion command or drawing code, so no image is delivered and the package declares no asset dependency.  Editorial formatting only, in addition: an AMS \texttt{amsart} preamble that restates the article's own declarations and macros used by the retained lines, namely the environments \texttt{as}, \texttt{lemma} on separate counters, together with the standard packages those lines need.  The retained environments print in this document as As 1, Lemma 1 in order of appearance, whereas the article prints the same items with the numbering of its own full document.  The complete native source of the article as distributed in the arXiv e-print for this exact version is bundled unchanged as \texttt{sources/203388/source0.tex}.
\begin{equation}\label{eq:sgm-ls-gap}
  g_\alpha(x_k+\lambda_k d_k)\le \mathcal T_k+\delta_1\lambda_k\inner{v_k}{d_k}-\delta_2\lambda_k^2\norm{d_k}^2 .
\end{equation}

\begin{as}\label{as:standing}
\begin{enumerate}[label=(\roman*),leftmargin=2em]
  \item $K\subseteq\R^n$ is nonempty, closed, and convex, and $\Sol\neq\emptyset$.
  \item The relevant gap function ($g_\alpha$ for VI/MVI, $g_\alpha^{\mathrm{EP}}$ for EP) is coercive, so
  its sublevel sets are bounded. Under strong monotonicity this is automatic, since the error bound makes
  the gap coercive.
  \item $F$ is continuous (resp.\ $G$ is continuous).
  \item The scaling matrices satisfy $\Dm_k\in\Dfam_\mu$ for all $k$.
\end{enumerate}
\end{as}

\begin{as}[Direction/sufficient-decrease condition]\label{as:dir}
There is $\kappa>0$ such that the certificate $v_k$ and direction $d_k$ satisfy
$\inner{v_k}{d_k}\le -\kappa\,\norm{d_k}^2$ for all $k$.
\end{as}

\begin{as}[V1 setting]\label{as:v1}
$F\in C^{1,1}$ is strongly monotone with modulus $m>0$ and $L$-Lipschitz. We take
$v_k=\nabla g_\alpha(x_k)$.
\end{as}

\begin{equation}\label{eq:dir-m}
  \inner{\nabla g_\alpha(x_k)}{d_k}\ \le\ -\,m\,\norm{d_k}^2 ,
\end{equation}

\begin{lemma}[Line search and sufficient decrease]\label{lem:descent}
Under Assumptions~\ref{as:standing} and~\ref{as:v1}:
\begin{enumerate}[label=(\roman*),leftmargin=2em]
  \item the line search terminates and
  $\lambda_k\ge\lambda_{\min}:=\min\{\kappa,\lambda_{\max},\sigma\bar\lambda\}>0$, where $\kappa$ is the
  direction constant of Assumption~\ref{as:dir} (here $\kappa=m$ by~\eqref{eq:dir-m}) and
  \[
  \bar\lambda=(1-\delta_1)\kappa/(L_g/2+\delta_2).
  \]
  \item $g_\alpha(x_k)\le\mathcal T_k$ for all $k$, and $\{\mathcal T_k\}$ is nonincreasing.
  \item $\mathcal T_k-\mathcal T_{k+1}\ge (1-\eta_{\max})\,\rho\,\norm{d_k}^2$ with
  $\rho=\delta_1\kappa\lambda_{\min}>0$.
\end{enumerate}
\end{lemma}

\begin{proof}
We argue by induction that $g_\alpha(x_k)\le\mathcal T_k$ for all $k$. This holds at $k=0$ because
$\mathcal T_0=g_\alpha(x_0)$. Assume it at index $k$, so $x_k\in\mathcal L_0$.
\begin{enumerate}[label=(\roman*),leftmargin=2em]
  \item For $0<\lambda\le\lambda_{\max}\le1$ the trial point $x_k+\lambda d_k$ lies on
  $[x_k,y_\alpha(x_k)]\subseteq\mathcal C_0$, so the descent lemma for the $L_g$-Lipschitz gradient on
  $\mathcal C_0$ gives
  \[
    g_\alpha(x_k+\lambda d_k)\ \le\ g_\alpha(x_k)+\lambda\inner{\nabla g_\alpha(x_k)}{d_k}
    +\tfrac{L_g}{2}\lambda^2\norm{d_k}^2 .
  \]
  With $g_\alpha(x_k)\le\mathcal T_k$ and $v_k=\nabla g_\alpha(x_k)$, the test \eqref{eq:sgm-ls-gap} holds
  whenever
  \[
    \big(\tfrac{L_g}{2}+\delta_2\big)\lambda\norm{d_k}^2\ \le\ (1-\delta_1)\big(-\inner{\nabla g_\alpha(x_k)}{d_k}\big) ,
  \]
  which by \eqref{eq:dir-m} is true for all $0<\lambda\le\bar\lambda$. Hence the backtracking terminates.
  Its first trial is $\min\{s_k,\lambda_{\max}\}\ge\min\{\kappa,\lambda_{\max}\}$ and the trials shrink by
  $\sigma\in(0,1)$. A trial fails only when it exceeds $\bar\lambda$, so the accepted step satisfies
  $\lambda_k\ge\min\{\kappa,\lambda_{\max},\sigma\bar\lambda\}=\lambda_{\min}$.

  \item The accepted step obeys \eqref{eq:sgm-ls-gap}, hence
  \[
    g_\alpha(x_{k+1})\ \le\ \mathcal T_k+\delta_1\lambda_k\inner{\nabla g_\alpha(x_k)}{d_k}
    -\delta_2\lambda_k^2\norm{d_k}^2\ \le\ \mathcal T_k .
  \]
  Therefore $\mathcal T_{k+1}=\eta_{k+1}\mathcal T_k+(1-\eta_{k+1})g_\alpha(x_{k+1})\le\mathcal T_k$, and
  \[
  \mathcal T_{k+1}-g_\alpha(x_{k+1})=\eta_{k+1}\big(\mathcal T_k-g_\alpha(x_{k+1})\big)\ge0.
  \]
  This closes the
  induction and shows $\{\mathcal T_k\}$ is nonincreasing.

  \item Using \eqref{eq:dir-m}, $\lambda_k\ge\lambda_{\min}$, and $\eta_{k+1}\le\eta_{\max}$,
  \[
  \begin{array}{lcl}
    \mathcal T_k-\mathcal T_{k+1}&=&(1-\eta_{k+1})\big(\mathcal T_k-g_\alpha(x_{k+1})\big) \\
    \ &\ge& \ (1-\eta_{\max})\big(-\delta_1\lambda_k\inner{\nabla g_\alpha(x_k)}{d_k}\big) \\
    \ &\ge&\ (1-\eta_{\max})\,\delta_1\kappa\lambda_{\min}\norm{d_k}^2=(1-\eta_{\max})\,\rho\,\norm{d_k}^2 .
  \end{array}
  \]
\end{enumerate}
\end{proof}

\end{document}
